\section{Groups of type \texorpdfstring{$D_4$}{D4}}\label{sec:d4}

For type $D_4$ there are no anisotropic nonarchimedean factors.  Our
objective is therefore to show that every finite quotient of the rational
group is trivial.  The argument first kills rational noncentral
involutions and then represents every power-subgroup coset by a product of
two such involutions.  The second step requires a maximal torus on which
inversion lifts rationally, not just geometrically.

\begin{theorem}\label{thm:d4}
Let $k$ be a number field and let $G$ be an absolutely almost simple simply
connected $k$-group of type $D_4$, including the trialitarian forms.  Every
noncentral abstract normal subgroup of $G(k)$ is $G(k)$.
\end{theorem}

The isotropic case is among the established cases recalled in
Section~\ref{sec:foundations}.  We may thus assume throughout this section
that $G$ is $k$-anisotropic.  Fix a positive integer $e$, put
$R=G(k)^e$, and write $q:G(k)\to\mathcal Q=G(k)/R$.  By
Proposition~\ref{prop:finite-defect}, this is a finite quotient and its
cosets have the approximation properties used below.  It will suffice to
prove that $\mathcal Q=1$.

\subsection{Root involutions and a quadratic splitting field}

We record explicitly the lattice calculations used both here and in
Section~\ref{sec:e6}.  With roots of squared length two, the coroot lattice
of the simply connected group is
\[
 Q_4=\{(x_1,x_2,x_3,x_4)\in\bbZ^4:x_1+x_2+x_3+x_4\equiv0\pmod2\},
 \qquad \Phi_4=\{\pm e_i\pm e_j:i\ne j\}.
\]
The two-torsion of a split maximal torus is $Q_4/2Q_4$.  Its four central
elements are the classes having all coordinates of the same parity.  The
other twelve classes are represented by the twelve pairs of opposite
roots.  These representatives are distinct: if two roots differ by an
element of $2Q_4$, their coordinate parities agree, and the defining
even-sum condition then leaves only equality or opposition.  The Weyl
group is transitive on the roots, so all noncentral involutions form one
geometric conjugacy class.

For $n=\alpha^\vee(-1)$, the roots of $C_G(n)$ are precisely the four
orthogonal axes represented, after a Weyl conjugacy, by
\[
 \alpha=e_1-e_2,\qquad \beta_1=e_1+e_2,
 \qquad\beta_2=e_3-e_4,\qquad\beta_3=e_3+e_4.
\]
The semisimple centralizer is connected, by the connected-centralizer
theorem for simply connected groups
\cite[Section~2.10, p.~53]{SteinbergRegular1965}. The lattice
spanned by these four roots has index two in $Q_4$, and
\[
 \alpha+\beta_1+\beta_2+\beta_3=2(e_1+e_3)\in2Q_4.
\]
Consequently the centralizer and its simply connected covering have the
geometric description
\begin{equation}\label{eq:d4-centralizer}
 C_G(n)_{\bar k}\simeq
 (\SL_2^4)/\langle(-I,-I,-I,-I)\rangle.
\end{equation}
The product of the four standard reflection representatives acts by $-1$
on the torus and has square one in this quotient.  Every other
representative of inversion has the same square: for $a$ in the torus
and an inverter $x$, one has $(ax)^2=a\,xax^{-1}x^2=x^2$.

\begin{lemma}\label{lem:d4-kill-involutions}
Every noncentral involution in $G(k)$ is killed by $q$.
\end{lemma}

\begin{proof}
Let $n\in G(k)$ be such an involution.  The root pair representing $n$ is
unique, so its axis $\alpha$ is Galois-stable.  The simply connected cover
of $C_G(n)$ is therefore a product of an inner form of $\SL_2$ over $k$
and restrictions of inner forms of $\SL_2$ over a cubic etale $k$-algebra
$E$.  Equivalently these factors are the norm-one groups of quaternion
algebras over $k$ and $E$.

There is a quadratic field $L'/k$ whose scalar extension embeds as a
maximal etale quadratic algebra in every one of these quaternion
algebras.  To see this, impose local conditions only at their ramified
places.  At a finite ramified place we require $L'_v$ to be a field and
to remain a field over each completion of $E$ where the corresponding
quaternion algebra is ramified.  An odd-degree completion cannot contain
$L'_v$.  Since $\dim_k E=3$, at most one completion of $E$ has degree two;
thus at most one quadratic local field is excluded.  There are other
quadratic extensions of a nonarchimedean local field.  At real ramified
places choose $L'_v=\bbC$.  Weak approximation on a defining square class,
with a further nonsquare condition if necessary, realizes all these
requirements.  The local invariant theorem for quaternion algebras then
shows that $L'$ splits each algebra after scalar extension, which is
equivalent to the asserted quadratic embedding
\cite{PRbook,Reiner,KMRT}.

The resulting maximal norm tori give a maximal torus of $G$.  Orient the
four root axes using their common quadratic field $L'$.  Galois then acts
by a common sign, together with permutations of $\beta_1,\beta_2,\beta_3$
fixing $\alpha$.  In particular the two roots
\[
 \alpha,\qquad
 \gamma=\frac{-\alpha+\beta_1+\beta_2+\beta_3}{2}
\]
span a Galois-stable subsystem of type $A_2$: both have squared length
two and $(\alpha,\gamma)=-1$.  The corresponding connected root subgroup
$J$ is defined over $k$ and is split over $L'$, because both roots and
coroots of this subsystem are fixed over $L'$.  The element $n$ is the
image of the rational element $\alpha^\vee(-1)$ in the simply connected
cover $\widetilde J$.

The group $\widetilde J$ has no anisotropic finite places.  In the inner
case its degree-three central simple algebra is split by a quadratic
extension, so its index divides both three and two.  In the outer case
the algebra over the quadratic center is split; a ternary Hermitian form
over a quadratic extension of a nonarchimedean local field is isotropic.
The established inner-type-$A$ result and Theorem~\ref{thm:outer-a} now
give (MP) for $\widetilde J$.  Lemma~\ref{lem:restriction-mp}, with an empty
anisotropic-place product, makes its image under $q$ trivial.  In
particular $q(n)=1$.
\end{proof}

\subsection{Local representatives of inversion}

The next two lemmas give the local conditions imposed on a representative
of an arbitrary coset of $R$.  At a prescribed place we can choose a
suitable open set freely.  At the other places we will use a uniform
condition on integral reduction.

\begin{lemma}\label{lem:d4-local-openings}
For every completion $F$ of $k$ there is a nonempty open set of regular
semisimple elements $u\in G(F)$ such that inversion on $C_G(u)$ has a
representative in $G(F)$.
\end{lemma}

\begin{proof}
We first produce a noncentral involution in $G(k)$, which can then be
localized at any $F$.  The unique geometric class of noncentral
involutions is a closed semisimple orbit preserved by Galois, so it
descends to a $k$-homogeneous space $X$ under $G$.  Its full geometric
stabilizer is the connected semisimple centralizer in
Equation~\eqref{eq:d4-centralizer}.  Borovoi's theorem for a homogeneous
space under a simply connected group with connected geometric
stabilizer having no torus quotient requires only archimedean local
points to conclude $X(k)\ne\varnothing$
\cite[Proposition~3.4(i)]{Borovoi}.

These points exist over $\bbC$.  Over $\bbR$, the compact rank of every
real form of $D_4$ is at least three, by the real orthogonal
classification, including the quaternionic orthogonal form
\cite{KMRT,PRbook}.  A compact torus of dimension at least three has at
least eight rational two-torsion points, whereas the center has order
four.  It therefore contains a noncentral involution.  Borovoi's theorem
now supplies $n\in G(k)$.  This argument applies to all forms, including
trialitarian ones; it imposes no finite-place restriction.

View $n$ in $G(F)$.  Geometrically it is conjugate
to an inversion representative by Equation~\eqref{eq:d4-centralizer} and
the preceding lattice calculation.  Hence the map
$g\mapsto gng^{-1}n$ takes some geometric values that are regular
semisimple: on a torus inverted by $n$, its values include all squares.
The regular locus is a nonempty Zariski open condition on $g$, and
$G(F)$ is Zariski dense.  We may therefore choose such a $g$ rationally.
Set $u=gng^{-1}n$ and $T=C_G(u)$.  The involution $n$ normalizes $T$ and
sends $u$ to $u^{-1}$.  Its Weyl action is inversion.  Indeed inversion
belongs to the $D_4$ Weyl group; composing the two Weyl actions would fix
$u$, whose Weyl stabilizer is trivial by connectedness of its semisimple
centralizer and regularity.  Finally the conjugation map
$G(F)\times T(F)\to G(F)$ is a submersion at $(1,u)$, since no root takes
the value one on $u$.  Thus nearby regular elements have conjugate tori,
and the required property holds on an open neighborhood of $u$.
\end{proof}

\begin{lemma}\label{lem:d4-one-wall}
Outside a fixed finite set of finite places, the following assertion
holds.  Let $u\in G(k_v)$ be regular semisimple and integral.  If the
semisimple part of its reduction lies on at most one pair of opposite
root walls, then inversion on $C_G(u)$ has a rational representative.
\end{lemma}

\begin{proof}
Write $p$ for the residue characteristic and exclude $p=2$ and the fixed
bad-model places.  The compact closure of the cyclic group generated by
$u$ has a prime-to-$p$ finite-order factor.  Its corresponding element
$s$ reduces to the semisimple part of the reduction of $u$.  This follows
also by decomposing the finite-order part and the pro-$p$ part of this
procyclic group, whose reduction kernel is pro-$p$.  In particular
$T=C_G(u)$ is contained in $C_G(s)$.

The centralizer $C_G(s)$ has an unramified reductive model with the roots
specified by the reduction of $s$.  One way to check this assertion is to
pass to a finite unramified extension splitting the reduced data.  A
prime-to-$p$ torsion lift is integrally conjugate to the corresponding
element of a split torus: at each successive congruence level the
conjugacy obstruction is cohomology of a finite cyclic group of order
prime to $p$ on a residue Lie algebra, and averaging makes this
cohomology vanish.  The resulting centralizer has precisely the root
groups on which $s$ has value one.  Descent gives the asserted unramified
model.  Equivalently one may use smoothness of fixed points of an
automorphism of invertible finite order together with the
connected-centralizer theorem.

If there are no root walls, this centralizer is the unramified torus
$T$.  Inversion is a rational point of its Weyl-group quotient.  It has a
representative over the residue field by Lang's theorem for the torus,
and that representative lifts through the smooth unramified normalizer.

If there is one root pair $\{\pm\alpha\}$, its element
$n_0=\alpha^\vee(-1)$ is Galois-invariant, even if Galois exchanges
$\alpha$ and $-\alpha$.  It is an integral noncentral involution central
in $C_G(s)$.  Thus $T\subset C_G(n_0)$.  The latter centralizer is again
unramified, and its simply connected covering is a product of
restrictions of split $\SL_2$'s over unramified extensions: a semisimple
unramified group is quasi-split, and a quasi-split group of type $A_1$ is
split.  The covering torus above $T$ is a product of maximal tori of
these $\SL_2$'s.

For each such torus, represent it as the norm-one group of a quadratic
etale algebra.  Conjugation of that algebra, acting on its underlying
two-dimensional vector space, is an inverter in $\GL_2$ of determinant
$-1$.  Multiply all these inverters, over all factors and embeddings,
by the same scalar $i$ with $i^2=-1$.  The resulting tuple lies in
$\SL_2^4$ and inverts the covering torus.  A Galois automorphism either
preserves $i$ or changes its sign simultaneously in all four factors.
The latter discrepancy is exactly the kernel in
Equation~\eqref{eq:d4-centralizer}.  The image of the tuple therefore
descends to a $k_v$-rational element of $C_G(n_0)$ inverting $T$.  Notice
that this argument permits the individual quadratic tori to be ramified.
\end{proof}

\subsection{A torus with no local-to-global obstruction}

Write $P_4=\Hom(Q_4,\bbZ)$ for the weight lattice and
$W=W(D_4)$ for its Weyl group.  The following calculation explains why
forcing the full Weyl group into a torus's Galois action is sufficient.

\begin{lemma}\label{lem:d4-cyclic-cohomology}
Let $\Gamma$ be a subgroup of $\Aut(\Phi_4)$ containing $W$, acting on
$P_4$.  Every class in $H^2(\Gamma,P_4)$ whose restriction to every cyclic
subgroup vanishes is zero.  Consequently, if a maximal $k$-torus $T$ of
a simply connected group of type $D_4$ has character action containing
$W$, then $\Sha^1(k,T)=0$.
\end{lemma}

\begin{proof}
Let $M=P_4$ and let $\xi\in H^2(\Gamma,M)$ restrict to zero on every
cyclic subgroup.  The element $c=-I\in W$ is central in $\Gamma$ and
acts by $-1$ on $M$.  Apply Lemma~\ref{lem:foundations-sign-extension}
to an extension representing $\xi$.  It gives a class
\[
 \beta\in H^1(\Gamma/\langle c\rangle,M/2M)
\]
whose vanishing is equivalent to the splitting of the extension.  We
view a representative of this class as a cocycle on $\Gamma$ and show
that it is a coboundary.

For a simple reflection $s_\alpha$, cyclic triviality and the last
assertion of Lemma~\ref{lem:foundations-sign-extension} give a lift
$d_{s_\alpha}\in M$ of $\beta(s_\alpha)$ with
$(1+s_\alpha)d_{s_\alpha}=0$.  Therefore
\[
 d_{s_\alpha}\in P_4\cap\mathbb Q\alpha=\mathbb Z\alpha.
\]
The last equality follows by pairing with a neighboring simple coroot.
The fundamental weights allow us to prescribe all the coefficients
$\langle b,\alpha^\vee\rangle$ modulo two independently.  Adding the
coboundary $(1-\gamma)b$ for a single $b\in P_4$ therefore makes
$\beta$ vanish on every simple reflection, and hence on $W$.

Moreover,
\begin{equation}\label{eq:d4-no-mod-two-invariants}
 (P_4/2P_4)^W=0.
\end{equation}
Indeed a root is primitive in $P_4$; for example its pairing with an
adjacent simple coroot is $-1$.  If $p$ is fixed modulo $2P_4$ by every
simple reflection, the formula
$p-s_\alpha p=\langle p,\alpha^\vee\rangle\alpha$ therefore makes
every simple-coroot pairing even.  Expansion in fundamental weights
gives $p\in2P_4$.  Since $W$ is normal in $\Gamma$, the cocycle identity
and $\beta|_W=0$ imply that $\beta(\gamma)$ is $W$-invariant for every
$\gamma$.  Equation~\eqref{eq:d4-no-mod-two-invariants} gives
$\beta=0$ on all of $\Gamma$.  The extension consequently splits by
Lemma~\ref{lem:foundations-sign-extension}, proving the first assertion.

For the arithmetic conclusion, take the minimal Galois splitting field
of $T$ and identify $X^*(T)$ with $P_4$.  The first assertion and
Lemma~\ref{lem:foundations-cyclic-detection} give
$\Sha^2_\omega(k,X^*(T))=0$.  The torus-duality conclusion of that
lemma then gives $\Sha^1(k,T)=0$, as required.
\end{proof}

We now construct the torus to which this lemma will apply.  Fix a
splitting field and a pinning of $G$ there.  At finitely many sufficiently
large finite places that split completely in this field, prescribe
unramified maximal tori with Frobenius in each conjugacy class of $W$.
Such tori exist by Weyl twisting over the residue field and smooth
lifting.  Each has rational inversion by the no-wall argument in
Lemma~\ref{lem:d4-one-wall}.  Choose regular elements in them and impose
small regular neighborhoods with the same local tori.

The labeling here matters.  At each chosen place fix a completion
embedding of the splitting field and use the fixed pinning.  After a
global identification of a new maximal torus with the pinned torus by
inner conjugation, these local identifications differ by Weyl elements,
not arbitrary diagram automorphisms.  The image of the Galois group
fixing the chosen splitting field thus meets every $W$-conjugacy class
inside $W$.  This image is all of $W$.  Indeed a proper subgroup of a
finite group cannot meet every conjugacy class: its transitive action on
cosets would have no element without a fixed point, contradicting the
average fixed-point count of one.  The full character action $\Gamma$
therefore contains $W$.

At each place, define the allowed set as the union of the elements that
are not regular semisimple and the regular semisimple elements
whose centralizer torus admits rational inversion. The auxiliary
regular neighborhoods just chosen force the resulting rational element
to be regular semisimple, hence regular semisimple at every completion.
For this element, membership in the allowed sets therefore gives
local rational inversion everywhere.

Apply Proposition~\ref{prop:power-approximation} to any prescribed coset
of $R$, with these auxiliary conditions and the openings from
Lemma~\ref{lem:d4-local-openings} at the other fixed exceptional places.
We spell out the two uniform conditions needed for that proposition.

At integral good places, exclude the closed set where the semisimple
part lies on two independent root walls.  In a maximal torus these
intersections have codimension at least two, as does their image under
the finite quotient by $W$.  The conjugacy-quotient map has fibers of
dimension $\dim G-\rank G$, so its inverse image still has codimension
at least two.  To see the fiber dimension directly, a fiber is a union
of conjugates of $su$, where $s$ is fixed semisimple and $u$ ranges over
the unipotent variety of $C_G(s)$; the latter has dimension
$\dim C_G(s)-\rank G$.  These descriptions and codimension bounds spread
to the integral models after excluding finitely many places
\cite[Theorem~6.11, p.~64]{SteinbergRegular1965}.
At every remaining integral place, Lemma~\ref{lem:d4-one-wall} applies
if the element is regular semisimple; otherwise it is in the allowed
set by definition.

At a single pole, the place splits in the fixed splitting field used in
Proposition~\ref{prop:power-approximation}.  Use a split eight-dimensional
vector representation.  Choose the cocharacter to take distinct values
on its eight weights, also after every relevant diagram translate.  This
requires avoidance of only finitely many hyperplanes in the cocharacter
lattice.  Order the exponents as $d_1<\cdots<d_8$.  In the single-pole
normal form $g_L\lambda(t)^e g_R$, require all initial principal minors
of $g_Rg_L$ in this order to be nonzero upon reduction.  These conditions
are nonempty open conditions: they hold at the identity.  For each pole
factor, at least one of the two disjoint complete basic parameter lists
lies entirely before or after that factor.  After cyclically combining
the side factors as $g_Rg_L$, this list occurs between fixed invertible
translates.  Its product map is dominant, so it meets the simultaneous
principal-minor open set.  This argument includes poles in the first or
last factor; no boundary factor needs parameters on both sides.  The
$j$th characteristic coefficient then has a unique
least-valuation term, of valuation
$e(d_1+\cdots+d_j)\ord_v(t)$.  The Newton polygon has eight segments of
horizontal length one with distinct integral slopes.  Hensel lifting
therefore splits the characteristic polynomial into distinct linear
factors over $k_v$.  The centralizer torus is split: the vector
representation has finite kernel, so its weights span the character
space over $\bbQ$, on which Galois now acts trivially.  Split tori have
rational inversion representatives by the lattice calculation above.
The finitely many extra model or denominator exceptions are treated by
the local openings, as provided by
Proposition~\ref{prop:power-approximation}.

We have proved that every coset of $R$ has a regular semisimple
representative $u$ for which $T=C_G(u)$ has everywhere locally rational
inversion and character action containing $W$.

\begin{proof}[Proof of Theorem~\ref{thm:d4}]
For the torus just constructed, the fiber over inversion in
$N_G(T)\to N_G(T)/T$ is a $T$-torsor.  It has a point over every
completion, and Lemma~\ref{lem:d4-cyclic-cohomology} makes its class zero.
Choose a rational representative $x$ of inversion.  The common-square
calculation following Equation~\eqref{eq:d4-centralizer} gives
$x^2=(ux)^2=1$.  Neither involution is central, because each induces
nontrivial inversion on the positive-dimensional torus $T$.  Thus
$u=(ux)x$ is killed by Lemma~\ref{lem:d4-kill-involutions}.  Every coset
of $R$ is trivial, so $G(k)/G(k)^e=1$.  By
Proposition~\ref{prop:finite-defect}, an arbitrary noncentral abstract
normal subgroup has finite index and contains such a power subgroup.
It is consequently all of $G(k)$.
\end{proof}

\section{Groups of type \texorpdfstring{$E_6$}{E6}}\label{sec:e6}

The last case reduces a generic rational element to a product from a
subgroup of type $D_4$ and the centralizer of a root involution.  That
centralizer has types $A_1$ and $A_5$.  The geometric product
decomposition alone is insufficient: we will identify its rational
descent obstruction as a torsor under a simply connected semisimple
group, and arrange the required local points.

\begin{theorem}\label{thm:e6}
Let $k$ be a number field and let $G$ be an absolutely almost simple simply
connected $k$-group of type $E_6$.  Every noncentral abstract normal
subgroup of $G(k)$ is $G(k)$.
\end{theorem}

Again $A=\varnothing$, and only the $k$-anisotropic case needs proof.
Fix $e\ge1$, set $R=G(k)^e$, and write
$q:G(k)\to\mathcal Q=G(k)/R$.  Our goal is to kill each coset of this
finite quotient.

\subsection{A rational root involution}

\begin{lemma}\label{lem:e6-root-involution}
The group $G(k)$ contains an involution geometrically conjugate to
$\alpha^\vee(-1)$ for a root $\alpha$.  Its centralizer $H$ is connected
semisimple of type $A_1A_5$, and its simply connected covering has
geometric form
\begin{equation}\label{eq:e6-centralizer-cover}
 \pi:\SL_2\times\SL_6\longrightarrow H_{\bar k},
 \qquad \ker\pi=\langle(-I_2,-I_6)\rangle.
\end{equation}
\end{lemma}

\begin{proof}
First work geometrically.  The roots centralizing $\alpha^\vee(-1)$ are
$\pm\alpha$ and the roots perpendicular to $\alpha$.  The latter form
$A_5$.  For example, the identity
$\sum_{\beta\in\Phi(E_6)}(\alpha,\beta)^2=4h^\vee=48$, with
$h^\vee=12$, shows that forty of the seventy roots other than
$\pm\alpha$ have pairing $\pm1$, leaving thirty perpendicular roots.
The resulting rank-five subsystem is $A_5$.  The determinant comparison
$\det(A_1)\det(A_5)/\det(E_6)=2\cdot6/3=4$ gives index two for its
coroot lattice together with $\alpha$ in the $E_6$ coroot lattice.
Each individual summand is saturated: the square of any nontrivial
saturation index would divide two or six, respectively.  The kernel is
therefore the diagonal subgroup of order two, proving
Equation~\eqref{eq:e6-centralizer-cover}.  Connectedness again follows
from the semisimple-centralizer theorem.

The geometric root-involution class is a closed semisimple orbit
preserved by Galois, so it descends to a $k$-homogeneous space $X$ under
$G$.  Its full geometric stabilizer is the connected semisimple group
just described.  The ambient group is simply connected, and the
stabilizer has no torus quotient.  Borovoi's theorem therefore gives a
$k$-point of $X$ as soon as $X$ has points at every archimedean place
\cite[Proposition~3.4(i)]{Borovoi}.  No finite-place existence assertion
is needed.  We verify the real places explicitly.

Over $\bbR$, take a real maximal torus and identify its cocharacter
lattice with $Q(E_6)$.  Complex conjugation is an involution on the
six-dimensional vector space $Q(E_6)/2Q(E_6)$, so its fixed space has
dimension at least three.  The quadratic form
\[
 \overline x\longmapsto (x,x)/2\pmod2
\]
on this space is nondegenerate, since $\det(E_6)=3$ is odd.  The
thirty-six root pairs give thirty-six distinct vectors of value one.
Indeed if $\alpha-\beta=2x$, then
$(x,x)=1-(\alpha,\beta)/2$; evenness of the root lattice permits this
only for $\beta=\pm\alpha$.  A nondegenerate quadratic form in dimension
six over $\bbF_2$ has either twenty-eight or thirty-six vectors of value
one.  Ours consequently has minus type and Witt index two, with the root
pairs exhausting its value-one vectors.  A three-dimensional fixed
subspace cannot be totally singular.  Its value-one vector is a real
two-torsion point in the required root-involution class.

Existence over $\bbC$ is immediate.  All the archimedean hypotheses of
Borovoi's theorem are now satisfied, giving a rational root involution
$n\in G(k)$ and completing the proof.
\end{proof}

Fix $n$ as in the lemma, put $H=C_G(n)$, and denote its simply connected
cover over $k$ by $\pi:\widetilde H\to H$.  Its almost simple factors
have types $A_1$ and $A_5$, for which (MP) is now known.

\subsection{The generic \texorpdfstring{$D_4$}{D4} subgroup}

\begin{lemma}\label{lem:e6-generic-d4}
There is a nonempty $k$-open subset $\mathcal U\subset G$ with the
following property.  For $u\in\mathcal U(k)$, set
$n'=unu^{-1}$ and $t=n'n$.  Then $t$ is regular semisimple, with maximal
torus $T=C_G(t)$, and the roots of $T$ that conjugation by $n$ sends to
their negatives generate a simply connected $k$-subgroup $D_*$ of type
$D_4$.  The group $D_*$ contains $n$ and $n'$ as geometrically conjugate
elements, and contains $t$ in its maximal torus $T_*=D_*\cap T$, on
which $n$ acts by inversion.
\end{lemma}

\begin{proof}
Work first over $\bar k$.  Number the $E_6$ diagram by the chain
$1,3,4,5,6$, with node $2$ attached to node $4$, and let $\delta$ be its
diagram involution, as in Figure~\ref{fig:e6-diagram}.  The four roots
\begin{figure}[tbp]
\centering
\begin{tikzpicture}[x=10mm,y=9mm,every label/.style={fill=white,inner sep=1pt}]
 \draw[gray,densely dashed] (2,-0.65)--(2,1.65);
 \draw (0,0)--(4,0) (2,0)--(2,1);
 \node[circle,draw,fill=white,inner sep=2pt,label=below:$1$] at (0,0) {};
 \node[circle,draw,fill=white,inner sep=2pt,label=below:$3$] at (1,0) {};
 \node[circle,draw,fill=white,inner sep=2pt,label=below:$4$] at (2,0) {};
 \node[circle,draw,fill=white,inner sep=2pt,label=below:$5$] at (3,0) {};
 \node[circle,draw,fill=white,inner sep=2pt,label=below:$6$] at (4,0) {};
 \node[circle,draw,fill=white,inner sep=2pt,label=above:$2$] at (2,1) {};
\end{tikzpicture}
\caption{The $E_6$ numbering used in
Equation~\eqref{eq:e6-fixed-roots}.  Reflection in the dashed axis is
$\delta$: it exchanges $1$ with $6$ and $3$ with $5$, and fixes $2,4$.}
\label{fig:e6-diagram}
\end{figure}
\begin{equation}\label{eq:e6-fixed-roots}
 \alpha_2,\quad \alpha_4,\quad
 \alpha_3+\alpha_4+\alpha_5,\quad
 \alpha_1+\alpha_3+\alpha_4+\alpha_5+\alpha_6
\end{equation}
have the $D_4$ Gram matrix: the first has pairing $-1$ with each of the
other three, which are mutually perpendicular.  They span the full
$\delta$-invariant sublattice of $Q(E_6)$.  Indeed successive differences
give $\alpha_3+\alpha_5$ and $\alpha_1+\alpha_6$, along with
$\alpha_2,\alpha_4$.  Their lattice is thus saturated.  Its
squared-length-two vectors are exactly its $D_4$ roots, so the
$\delta$-fixed roots form this subsystem.  The associated subgroup $D$
is simply connected, since its coroot lattice is saturated in that of
$G$.

Let $T_0$ be the ambient maximal torus and $S$ the rank-four maximal
torus of $D$.  There is no root perpendicular to $S$.  Such a root
would belong to the anti-invariant space of $\delta$ and satisfy
$\delta\alpha=-\alpha$, impossible because a diagram automorphism
preserves positive roots.  Thus a nonempty open subset $S_{\mathrm{reg}}$
consists of elements regular in $G$, with centralizer $T_0$.

Inversion in $D$ acts on the full root space as $-1$ on the
$\delta$-invariant space and as $1$ on its perpendicular complement.
It therefore represents the longest Weyl element $w_0=-\delta$ of
$E_6$.  By the $D_4$ calculation it has a representative of order two,
which is conjugate inside $D$ to a root involution.  We may conjugate our
chosen $n$ geometrically so that it is this representative.  For every
$s\in S$, the element $sn$ is conjugate to $n$ inside $D$: choose
$r\in S$ with $r^2=s$ and use $rnr^{-1}=sn$.

Consider the morphism into the $G$-conjugacy class of $n$ given by
\[
 H\times S_{\mathrm{reg}}\longrightarrow \operatorname{Cl}_G(n),
 \qquad (h,s)\longmapsto h(sn)h^{-1}.
\]
Its fibers have dimension two.  In fact equality with the value at
$(h_0,s_0)$ implies, on multiplying by $n$, that
$a=h_0^{-1}h\in H$ satisfies $asa^{-1}=s_0$.  Regularity forces
$a\in H\cap N_G(T_0)$, and $a$ determines $s$.  Conversely each such
$a$ preserves $S=\im(1-n:T_0\to T_0)$, since it commutes with $n$;
this holds for every component of $H\cap N_G(T_0)$.  It also preserves
$S_{\mathrm{reg}}$, so it gives the corresponding input in the same
fiber.  The identity
component of $H\cap N_G(T_0)$ is the fixed subtorus of $w_0=-\delta$,
of dimension two.  As $\dim H=3+35=38$, the image has dimension
\[
 38+4-2=40=78-38=\dim\operatorname{Cl}_G(n).
\]
It therefore contains an open subset of the class.

For every point in this image, the asserted subgroup is $hDh^{-1}$ and
its torus is $hSh^{-1}$: multiplication by $n$ gives the regular element
$hsh^{-1}$.  Intrinsically, these are recovered from the roots on which
$n$ acts as minus one.  This intrinsic description is Galois-invariant.
The resulting property holds on a Galois-stable dense constructible
subset, which contains a nonempty $k$-open subset.  Pulling it back under
$u\mapsto unu^{-1}$ gives $\mathcal U$.  The same intrinsic description
then defines $D_*$ and $T_*$ over $k$, with all the claimed properties.
\end{proof}

\subsection{The full pullback and the factorization torsor}

We next justify the rational descent of the preceding geometric
factorization.  It is important here to identify the entire pullback
subgroup, not just its identity component.

\begin{lemma}\label{lem:e6-pullback}
For $u\in\mathcal U(k)$, let $D_*$ be as in
Lemma~\ref{lem:e6-generic-d4}, and put
\[
 J=C_{D_*}(n)=D_*\cap H,
 \qquad \widetilde J=\pi^{-1}(J)\subset\widetilde H.
\]
Then $\widetilde J$ is connected semisimple simply connected.  The
variety
\begin{equation}\label{eq:e6-factorization-torsor}
 Y_u=\{(w,\widetilde q)\in D_*\times\widetilde H:
                    w\pi(\widetilde q)=u\}
\end{equation}
is a torsor under $\widetilde J$ and has points over every
nonarchimedean completion of $k$.
\end{lemma}

\begin{proof}
Over $\bar k$, conjugate $n$ inside $D_*$ to the root representative.
The four rank-one factors covering $J$ then map to
$\widetilde H_{\bar k}=\SL_2\times\SL_6$ by
\begin{equation}\label{eq:e6-four-block-map}
 (a,b,c,d)\longmapsto\bigl(a,\diag(b,c,d)\bigr).
\end{equation}
The distinguished root factor maps to the $A_1$ factor, and the three
perpendicular factors map to three disjoint two-dimensional blocks of
$A_5$.  This is an injective homomorphism.  The simultaneous minus one
maps to $(-I_2,-I_6)$, the full kernel of $\pi$, and the resulting
quotient is exactly $J$, by Equation~\eqref{eq:d4-centralizer}.  Hence
the image of Equation~\eqref{eq:e6-four-block-map} contains the kernel
of $\pi$ and surjects onto $J$.  It is therefore the full inverse image
$\pi^{-1}(J)$, not merely its connected component.  This identifies that
inverse image geometrically with $\SL_2^4$; descent proves that
$\widetilde J$ is connected semisimple simply connected over $k$.

The variety $Y_u$ is geometrically nonempty, since $n'$ is conjugate to
$n$ inside $D_*$.  Indeed if $wnw^{-1}=n'$, then $w^{-1}u\in H$, and
this element lifts through $\pi$ over $\bar k$.  The right action
\[
 (w,\widetilde q)\cdot y
       =\bigl(w\pi(y),y^{-1}\widetilde q\bigr),
       \qquad y\in\widetilde J,
\]
is simply transitive on geometric points.  For two factorizations, the
ratio of their $D_*$ factors lies in $D_*\cap H=J$, and the ratio of
their covering factors is its unique specified lift in
$\widetilde J$.  Thus Equation~\eqref{eq:e6-factorization-torsor} is a
$\widetilde J$-torsor.  Local $H^1$-vanishing for simply connected
semisimple groups gives its points at every nonarchimedean place
\cite{PRbook}.
\end{proof}

To use this torsor, we must choose local factorizations that are small
in the $\widetilde H$ factor.  The next elementary observation ensures
that the square root used for this purpose remains in the varying torus
$T_*$.

Fix a faithful representation of $G$.  At any completion, sufficiently
near the identity, the matrix analytic square root $t\mapsto t^{1/2}$
is defined and commutes with conjugation.  For all the tori $T_*$ in
Lemma~\ref{lem:e6-generic-d4}, one can choose this neighborhood uniformly
so that
\begin{equation}\label{eq:e6-uniform-root}
                  t\in T_*\quad\Longrightarrow\quad t^{1/2}\in T_*.
\end{equation}
Indeed these tori form one geometric conjugacy class of subtori, and
$t$ is regular semisimple, so its image in the fixed representation is
diagonalizable.  In such a diagonalization their character relations
are a fixed finitely generated lattice of monomial relations, up to a
permutation of the eigenvalues.  The square root is defined by a
convergent matrix power series and commutes with conjugation.  Its
eigenvalues are uniformly near one when the original matrix is near
one: a bound on the matrix norm bounds the absolute values of all its
eigenvalues, without any bound on a diagonalizing matrix.  Every
defining monomial evaluated on the root eigenvalues has square one,
because $t\in T_*$.  After shrinking a single neighborhood using the
finite list of relations, each such monomial must be one, proving
Equation~\eqref{eq:e6-uniform-root}.  This argument applies also at the
places above two, by choosing a smaller analytic neighborhood.

\begin{proof}[Proof of Theorem~\ref{thm:e6}]
The restriction of $q\circ\pi$ to $\widetilde H(k)$ is controlled by
(MP) for its $A_1$ and $A_5$ factors, using Theorem~\ref{thm:outer-a} in
the outer case.  By Lemma~\ref{lem:restriction-mp}, it extends to a
continuous homomorphism on their finite products of anisotropic finite
local groups.  Choose a finite set $B$ containing these places and all
archimedean places.  At each finite place in $B$, a sufficiently small
identity neighborhood in $\widetilde H(k_v)$ is killed by this local
homomorphism.

Take any coset of $R$.  The approximation conclusion of
Proposition~\ref{prop:finite-defect} gives a representative
$u\in\mathcal U(k)$ arbitrarily close to one at all places of $B$.
Use the notation of Lemma~\ref{lem:e6-generic-d4}.  For $v\in B$, the
element $t=unu^{-1}n$ is close to one.  Set
$w_v=t^{1/2}\in T_*(k_v)$, using
Equation~\eqref{eq:e6-uniform-root}.  Since $n$ inverts $T_*$, we have
\[
 w_v n w_v^{-1}=w_v^2n=tn=n'.
\]
Thus $w_v^{-1}u\in H(k_v)$ is close to one.  The central isogeny
$\pi$ is an analytic local isomorphism at one, so this element lifts to
$\widetilde q_v\in\widetilde H(k_v)$ close to one.  At the finite
places in $B$ the lift can be placed in the chosen kernel neighborhoods.
At the archimedean places this construction supplies the local point of
$Y_u$ that is needed for the Hasse principle.

Lemma~\ref{lem:e6-pullback} supplies points of $Y_u$ at all other finite
places.  The Hasse principle for torsors under simply connected
semisimple groups gives a rational point.  After trivializing the torsor
by this point, weak approximation for the same groups gives a rational
point $(w,\widetilde q)$ approximating the specified local points at $B$
\cite{PRbook}.  Consequently
\[
                         u=w\pi(\widetilde q).
\]
Theorem~\ref{thm:d4} kills $w\in D_*(k)$ under $q$, since a group of type
$D_4$ has no anisotropic finite places.  The local kernel conditions and
Lemma~\ref{lem:restriction-mp} kill $\pi(\widetilde q)$.  Thus $q(u)=1$.
Every coset of $R$ has been killed, so $G(k)/G(k)^e=1$.  The finite-index
conclusion of Proposition~\ref{prop:finite-defect} now gives the stated
assertion for every noncentral abstract normal subgroup.
\end{proof}

\subsection{Completion of the normal-subgroup theorem}

\begin{proof}[Proof of Theorem~\ref{thm:main}]
The established cases recalled in Section~\ref{sec:foundations},
Theorem~\ref{thm:outer-a}, Theorem~\ref{thm:d4}, and
Theorem~\ref{thm:e6} exhaust the absolutely almost simple simply
connected groups over number fields.  Their proofs give, in every case,
the vanishing of the finite defect for the abstract power subgroup
$R=G(k)^e$:
\[
 R=\delta^{-1}(P_0),\qquad
 P_0=\overline{\delta(R)}\triangleleft\prod_{v\in A}G(k_v),
 \qquad P_0\ \text{open}.
\]
For the established cases this equality follows directly from (MP);
for the remaining cases it is the conclusion of the preceding
arguments.

Let now $N\triangleleft G(k)$ be any abstract normal subgroup not
contained in the center.  Proposition~\ref{prop:finite-defect} gives
$[G(k):N]<\infty$.  Choose $e$ to be the exponent of the finite group
$G(k)/N$, so that $R\subset N$.  Density of the diagonal, and openness
of $P_0$, identify
\[
 G(k)/R\simeq
 \left(\prod_{v\in A}G(k_v)\right)/P_0.
\]
The normal subgroup $N/R$ therefore corresponds to a normal subgroup of
this finite local quotient.  Its full inverse image $W$ in
$\prod_{v\in A}G(k_v)$ is open and normal, and
$N=\delta^{-1}(W)$.

If $A$ is empty, the local product, its quotient, and $W$ are all the
one-element group.  The same argument gives $R=G(k)$ and hence
$N=G(k)$.  No finite-generation or closedness hypothesis on $N$ has
entered the argument.
\end{proof}
